%% =================================================================
%% Kingma, Ba. "Adam: A Method for Stochastic Optimization."
%% ICLR 2015. arXiv:1412.6980v9.
%%
%% Reconstruction of page 10 (printed p. 10) as study notes.
%% Generated by: reproduce-all-pages-basic-tex (batch 1-15)
%%
%% Structure: end of Sec 8 Conclusion + Sec 9 ACKNOWLEDGMENTS + start
%% of REFERENCES. No equations.
%% New content type: printed bibliography as page content (vs.
%% citation markers). Using a hanging-indent list.
%% =================================================================

\documentclass[11pt]{article}

\usepackage[margin=1in]{geometry}
\usepackage{amsmath, amssymb}
\usepackage{mathrsfs}
\usepackage{bm}
\usepackage{xcolor}
\usepackage{fancyhdr}

\pagestyle{fancy}
\fancyhf{}
\lhead{\small Published as a conference paper at ICLR 2015}
\rhead{}
\cfoot{\small 10 $\to$ \arabic{page}}
\renewcommand{\headrulewidth}{0.4pt}

\setcounter{page}{1}
\setcounter{equation}{12}

\begin{document}

% ---------- End of Conclusion ----------
\noindent
large datasets and/or high-dimensional parameter spaces. The
method combines the advantages of two recently popular
optimization methods: the ability of AdaGrad to deal with sparse
gradients, and the ability of RMSProp to deal with non-stationary
objectives. The method is straightforward to implement and
requires little memory. The experiments confirm the analysis on
the rate of convergence in convex problems. Overall, we found
Adam to be robust and well-suited to a wide range of non-convex
optimization problems in the field machine learning.

\vspace{0.5em}

\section*{9 \textsc{Acknowledgments}}

This paper would probably not have existed without the support
of Google Deepmind. We would like to give special thanks to Ivo
Danihelka, and Tom Schaul for coining the name Adam. Thanks to
Kai Fan from Duke University for spotting an error in the
original AdaMax derivation. Experiments in this work were partly
carried out on the Dutch national e-infrastructure with the
support of SURF Foundation. Diederik Kingma is supported by the
Google European Doctorate Fellowship in Deep Learning.

\vspace{0.5em}

\section*{\textsc{References}}

\begin{list}{}{%
  \setlength{\leftmargin}{1.5em}%
  \setlength{\itemindent}{-1.5em}%
  \setlength{\itemsep}{0.5em}%
  \setlength{\parsep}{0pt}%
}

\item Amari, Shun-Ichi. Natural gradient works efficiently in
learning. \emph{Neural computation}, 10(2):251--276, 1998.

\item Deng, Li, Li, Jinyu, Huang, Jui-Ting, Yao, Kaisheng, Yu,
Dong, Seide, Frank, Seltzer, Michael, Zweig, Geoff, He, Xiaodong,
Williams, Jason, et al. Recent advances in deep learning for
speech research at microsoft. \emph{ICASSP 2013}, 2013.

\item Duchi, John, Hazan, Elad, and Singer, Yoram. Adaptive
subgradient methods for online learning and stochastic
optimization. \emph{The Journal of Machine Learning Research},
12:2121--2159, 2011.

\item Graves, Alex. Generating sequences with recurrent neural
networks. \emph{arXiv preprint arXiv:1308.0850}, 2013.

\item Graves, Alex, Mohamed, Abdel-rahman, and Hinton, Geoffrey.
Speech recognition with deep recurrent neural networks. In
\emph{Acoustics, Speech and Signal Processing (ICASSP), 2013
IEEE International Conference on}, pp.~6645--6649. IEEE, 2013.

\item Hinton, G.E. and Salakhutdinov, R.R. Reducing the
dimensionality of data with neural networks. \emph{Science},
313(5786):504--507, 2006.

\item Hinton, Geoffrey, Deng, Li, Yu, Dong, Dahl, George E,
Mohamed, Abdel-rahman, Jaitly, Navdeep, Senior, Andrew,
Vanhoucke, Vincent, Nguyen, Patrick, Sainath, Tara N, et al.
Deep neural networks for acoustic modeling in speech
recognition: The shared views of four research groups.
\emph{Signal Processing Magazine, IEEE}, 29(6):82--97, 2012a.

\item Hinton, Geoffrey E, Srivastava, Nitish, Krizhevsky, Alex,
Sutskever, Ilya, and Salakhutdinov, Ruslan R. Improving neural
networks by preventing co-adaptation of feature detectors.
\emph{arXiv preprint arXiv:1207.0580}, 2012b.

\item Kingma, Diederik P and Welling, Max. Auto-Encoding
Variational Bayes. In \emph{The 2nd International Conference on
Learning Representations (ICLR)}, 2013.

\item Krizhevsky, Alex, Sutskever, Ilya, and Hinton, Geoffrey E.
Imagenet classification with deep convolutional neural networks.
In \emph{Advances in neural information processing systems},
pp.~1097--1105, 2012.

\item Maas, Andrew L, Daly, Raymond E, Pham, Peter T, Huang,
Dan, Ng, Andrew Y, and Potts, Christopher. Learning word vectors
for sentiment analysis. In \emph{Proceedings of the 49th Annual
Meeting of the Association for Computational Linguistics: Human
Language Technologies-Volume 1}, pp.~142--150. Association for
Computational Linguistics, 2011.

\item Moulines, Eric and Bach, Francis R. Non-asymptotic
analysis of stochastic approximation algorithms for machine
learning. In \emph{Advances in Neural Information Processing
Systems}, pp.~451--459, 2011.

\item Pascanu, Razvan and Bengio, Yoshua. Revisiting natural
gradient for deep networks. \emph{arXiv preprint
arXiv:1301.3584}, 2013.

\item Polyak, Boris T and Juditsky, Anatoli B. Acceleration of
stochastic approximation by averaging. \emph{SIAM Journal on
Control and Optimization}, 30(4):838--855, 1992.

\end{list}

% [page ends here — continues on page 11]

\end{document}
